\documentclass[11pt,a4paper]{article}

\usepackage[utf8]{inputenc}

\usepackage[T1]{fontenc}

\usepackage[polish]{babel}

\usepackage{lmodern}

\usepackage{geometry}

\geometry{a4paper,margin=27mm}

\usepackage{microtype}

\usepackage[hidelinks]{hyperref}

% Portfolio: MSC-ADHD-P04; source: prospective ARC-XR research scope.

\title{Transfer regulacji do ucieleśnionej interakcji i AR}

\author{Mateusz Pomianek}

\date{Wersja 0.3.0, 21 listopada 2026}

\begin{document}

\maketitle

\begin{abstract}

Przedmiotem opracowania jest Transfer regulacji do ucieleśnionej interakcji i AR. Czy umiejętność regulacji nabyta w grze ekranowej przenosi się na zadania ucieleśnione i działania w AR? Plan zakłada porównanie poprawności działania i przebiegu regulacji w grze ekranowej, zadaniu fizycznym oraz kontrolowanym zadaniu AR o wyrównanej trudności. Głównym punktem końcowym jest różnica jakości regulacji i poprawności działania w warunku transferowym względem bazowego. Opis ma charakter wstępny i nie przedstawia niezweryfikowanych wyników.

\end{abstract}

\section{Wprowadzenie}

Rozwój adaptacyjnych systemów XR stwarza możliwość badania interakcji człowieka z informacją przestrzenną, ale samo zastosowanie AR, VR lub modeli sztucznej inteligencji nie przesądza o przewadze poznawczej czy ergonomicznej. W rozpatrywanym zagadnieniu, Transfer regulacji do ucieleśnionej interakcji i AR, kluczowe jest rozdzielenie właściwości systemu od zmienności zachowania użytkownika. Porównanie należy zaprojektować tak, aby mierzyć jednoznaczne zjawisko, a nie tylko ogólną ocenę technologii.

\section{Cel pracy i pytanie badawcze}

Czy umiejętność regulacji nabyta w grze ekranowej przenosi się na zadania ucieleśnione i działania w AR? W zakresie ARC-XR chodzi o porównywalny kontrast transferu pomiędzy interakcją ekranową, ruchem i AR; materiał badania źródłowego pozostaje w repozytorium prowadzącym projekt.

\section{Zakres i odrębny wkład naukowy}
Transfer wyuczonej strategii regulacji z interfejsu ekranowego do zadania ucieleśnionego i AR. Dane i podstawowa odpowiedzialność naukowa pozostają w projekcie neuropsychologicznym, nie w odrębnym badaniu klinicznym ARC-XR. Transfer strategii przez zmianę środowiska, nie skuteczność samej gry lub nieuprawnione rozpoznawanie ADHD.

\section{Hipotezy badawcze}
\paragraph{H1.} Trening regulacji z jawną strategią poprawia wynik niećwiczonego zadania AR względem kontrolnego treningu o wyrównanej długości.
\paragraph{H2.} Transfer jest słabszy przy większej różnicy wymagań motorycznych i poznawczych między zadaniem treningowym a transferowym.
\paragraph{H3.} Korzyść w zadaniu ucieleśnionym utrzymuje się po odroczeniu, jeśli zachowana jest porównywalna trudność i równoległe wersje zadania.

\section{Projekt badania i plan eksperymentów}
\subsection{A. Operacjonalizacja i kwalifikacja danych}
Transfer wyuczonej strategii regulacji z interfejsu ekranowego do zadania ucieleśnionego i AR. Dane i podstawowa odpowiedzialność naukowa pozostają w projekcie neuropsychologicznym, nie w odrębnym badaniu klinicznym ARC-XR. Jednostka wnioskowania: uczestnik z powtórzonymi zadaniami transferowymi w trzech środowiskach. Punkt końcowy i sposób oceny będą zamrożone przed próbą potwierdzającą. Należy udokumentować pochodzenie obserwacji, czas akwizycji, identyfikatory sesji, harmonogram oraz legalność użycia materiałów.

\subsection{B. Eksperyment porównawczy}
Przygotować neutralne zadania ekranowe, fizyczne i AR o porównywalnej strukturze; randomizować trening, kontrolować ekspozycję. Estymand pierwszorzędowy: różnica poprawności działania w niećwiczonym zadaniu AR między warunkiem treningu regulacji i dopasowaną kontrolą ekspozycji. Przypisanie warunków, mianowniki i kryteria wyłączeń nie mogą być dobierane po przejrzeniu wyników.

\subsection{C. Kontrole i analiza odporności}
Nie wyprowadzać wniosków klinicznych ADHD ani o populacji specjalnej z niezweryfikowanej wygodnej próby; odrębnie zatwierdzić etykę. Zrównoważyć trudność, zmęczenie, efekty kolejności i powtórzenie, a materiał źródłowy powiązać po kanonicznym SHA. Każdą odmowę systemu, niekompletną próbę i brak sygnału należy raportować osobno, a nie po cichu usuwać z próby.

\subsection{D. Walidacja na niezależnym materiale}
Ocenę należy wykonać na uprzednio odłożonych scenach, sesjach lub uczestnikach. Model mieszany poprawności z interakcją trening × środowisko i uczestnikiem jako efektem losowym; porównanie odroczone pod kontrolą ekspozycji. Model i próg zatrzymać przed analizą końcową; rejestrować również brak przewagi i działania niepożądane.

\section{Estymand, analiza statystyczna i decyzje}
Główny estymand: różnica poprawności działania w niećwiczonym zadaniu AR między warunkiem treningu regulacji i dopasowaną kontrolą ekspozycji. Jednostką powtarzalności jest uczestnik z powtórzonymi zadaniami transferowymi w trzech środowiskach. Model mieszany poprawności z interakcją trening × środowisko i uczestnikiem jako efektem losowym; porównanie odroczone pod kontrolą ekspozycji. H1 jest kandydatem na hipotezę potwierdzającą. H2 i H3 pozostają eksploracyjne, dopóki protokół nie określi ich hierarchii i kontroli wielokrotnych porównań. Minimalny efekt istotny praktycznie, liczebność prób, margines nie gorszości i progi jakości wyznaczyć na podstawie pilotażu i analizy precyzji, nie zgadywać wartości.

\section{Stan badań i granica nowości}
\begin{itemize}
\item XR personalization methods exist; a claim needs a defined multi-session or transfer contrast. \cite{xrprior1}.
\item Learning-load effects do not imply delayed knowledge retention or skill transfer. \cite{xrprior2}.
\end{itemize}
Transfer strategii przez zmianę środowiska, nie skuteczność samej gry lub nieuprawnione rozpoznawanie ADHD. Samo zastosowanie VLM, EEG, okulometrii lub adaptacyjnego HUD nie stanowi wystarczającej nowości międzynarodowej.

\section{Warunki bezpieczeństwa, sprzętu i zgodności}
Opis stanowi prospektywny plan, a nie prerejestrację ani zatwierdzoną procedurę. Przed zbieraniem danych od ludzi niezbędna jest właściwa weryfikacja etyczna i zgoda. Standardowy Quest 3 nie ma wbudowanego eye trackingu, a Tobii Eye Tracker 5 działa na stanowisku ekranowym. Odstępstwa, obciążenia i ograniczenia czujników muszą być ujawnione. Nie publikować surowych danych biometrycznych ani nagrań uczestników.

\section{Materiał i metody}

Plan zakłada porównanie poprawności działania i przebiegu regulacji w grze ekranowej, zadaniu fizycznym oraz kontrolowanym zadaniu AR o wyrównanej trudności. Wymagane jest zdefiniowanie jednostki analizy, zasad doboru przypadków, wersji scenariusza i źródeł pomiaru. Rozkład warunków oraz reguły wykluczenia prób należy określić przed pomiarem. W badaniach z udziałem ludzi niezbędna jest uprzednia właściwa ocena etyczna i świadoma zgoda. W testach wyłącznie symulacyjnych należy oddzielić prawdę symulacyjną od obserwacji empirycznej.

\subsection{Wynik główny i kryteria jakości}

Różnica jakości regulacji i poprawności działania w warunku transferowym względem bazowego. Pomiar powinien zachowywać jawne jednostki, znaczniki czasu i reguły oceny poprawności. Przy modelach predykcyjnych stosować rozłączne partycje oceny; przy eksperymentach interwencyjnych oddzielić skuteczność predykcji od przyczynowego wpływu interfejsu.

\section{Plan analizy}

Model mieszany uwzględniający uczestnika, typ środowiska, kolejność i stopień doświadczenia; osobno ocenić odroczony transfer. W wynikach należy raportować wielkość efektu lub błąd estymacji z przedziałem niepewności, liczbą obserwacji i odsetkiem odrzuconych prób. Analizy dodatkowe wymagają jasnego oznaczenia jako eksploracyjne. Nie wybierać progów decyzyjnych ani modelu na podstawie obserwacji ze zbioru testowego.

\section{Wyniki}

W obecnej wersji nie raportuje się wyników pomiarów, tabel statystycznych ani parametrów wyznaczonych z nieudokumentowanych danych. Nie stwierdza się w tym opracowaniu, że eksperyment został wykonany lub zgłoszony do czasopisma.

\section{Dyskusja}

Interpretacja przyszłych wyników musi uwzględnić alternatywne mechanizmy i zależność wniosków od jakości danych, warunków zadania oraz ograniczeń sprzętu. Wartość potencjalnego rozwiązania zależy nie tylko od trafności modelu, lecz także od jego opóźnienia, stabilności i wpływu na zachowanie człowieka. Wyniki testów syntetycznych nie mogą być automatycznie uogólniane na populację uczestników.

\section{Ograniczenia badania}

Wynik może zależeć od nieporównywalności zadań, zmęczenia oraz uczenia specyficznego dla środowiska. Granice nowości i trafności ekologicznej wymagają niezależnej weryfikacji. Żadna część szkicu nie dowodzi przeprowadzenia zbierania danych, uzyskania zgód etycznych ani gotowości manuskryptu do zgłoszenia.

\section{Wnioski}

Przedstawiono problem, kontrast i sposób jego planowanej oceny. Rozstrzygnięcie postawionego pytania wymaga przygotowania protokołu, zweryfikowanych źródeł oraz poprawnie udokumentowanych wyników. Bez tych materiałów nie można sformułować konkluzji empirycznych.

\begin{thebibliography}{9}
\bibitem{xrprior1} Adapting to the User: A Systematic Review of Personalized Interaction in VR (2025), arXiv:2510.13123. \url{https://arxiv.org/abs/2510.13123}.
\bibitem{xrprior2} Can augmented reality technology reduce learners' cognitive load? A meta-analysis (2026), Smart Learning Environments. \url{https://doi.org/10.1186/s40561-025-00429-7}.
\end{thebibliography}
\end{document}

